\documentclass[12pt, a4paper]{article}

% ── Packages ──────────────────────────────────────────────────────────────────
\usepackage[utf8]{inputenc}
\usepackage[T1]{fontenc}
\usepackage{lmodern}
\usepackage[margin=1in]{geometry}
\usepackage{amsmath, amssymb, amsthm}
\usepackage{mathtools}
\usepackage{bm}
\usepackage{algorithm}
\usepackage{algpseudocode}
\usepackage{listings}
\usepackage{xcolor}
\usepackage{booktabs}
\usepackage{longtable}
\usepackage{enumitem}
\usepackage{hyperref}
\usepackage{fancyhdr}
\usepackage{titlesec}
\usepackage{parskip}
\usepackage{caption}
\usepackage{subcaption}
\usepackage{float}

% ── Black and white only ──────────────────────────────────────────────────────
\hypersetup{
    colorlinks=false,
    pdfborder={0 0 0},
}

% ── Code listing style (black and white) ─────────────────────────────────────
\lstdefinestyle{bwpython}{
    language=Python,
    basicstyle=\ttfamily\small,
    keywordstyle=\bfseries,
    commentstyle=\itshape,
    stringstyle=\ttfamily,
    numbers=left,
    numberstyle=\tiny,
    numbersep=8pt,
    frame=single,
    breaklines=true,
    breakatwhitespace=true,
    showstringspaces=false,
    tabsize=4,
    captionpos=b,
}
\lstset{style=bwpython}

% ── Theorem environments ─────────────────────────────────────────────────────
\newtheorem{definition}{Definition}[section]
\newtheorem{theorem}{Theorem}[section]
\newtheorem{remark}{Remark}[section]
\newtheorem{example}{Example}[section]

% ── Header / Footer ──────────────────────────────────────────────────────────
\pagestyle{fancy}
\fancyhf{}
\fancyhead[L]{\small YQuantum 2026 --- Hartford Portfolio QUBO}
\fancyhead[R]{\small\thepage}
\renewcommand{\headrulewidth}{0.4pt}

% ── Title ─────────────────────────────────────────────────────────────────────
\title{
    \textbf{YQuantum 2026: Hartford Portfolio QUBO} \\[0.5em]
    \large A Complete Technical Reference \\[0.3em]
    \large Theory, Implementation, Results, and Lessons Learned
}
\author{YQuantum Hackathon Team}
\date{April 2026}

\begin{document}
\maketitle
\thispagestyle{empty}

\begin{abstract}
This document is a self-contained reference for our YQuantum 2026 hackathon project:
an 8-qubit, insurer-aware portfolio optimization pipeline built for the
Capgemini $\times$ The Hartford challenge. We cover every piece of theory our
code depends on, explain every design decision, walk through each source file,
present our results, and discuss what we tried that did not work and why. The
intended audience is ourselves and any judge or teammate who wants to understand
the full depth of the project without reading the source code.
\end{abstract}

\newpage
\tableofcontents
\newpage

% ══════════════════════════════════════════════════════════════════════════════
\section{Introduction and Problem Statement}
% ══════════════════════════════════════════════════════════════════════════════

\subsection{The Challenge}

The Hartford is a major U.S.\ insurance company. Insurance companies hold large
investment portfolios to back their liabilities. The question we are asked to
solve is:

\begin{quote}
\textit{Given a universe of assets organized into sectors, which combination of
sectors should an insurer select to form a portfolio that balances return, risk,
regulatory capital requirements, and liquidity?}
\end{quote}

This is a \textbf{combinatorial optimization problem}: we must choose a subset
of sectors (a discrete decision), and the quality of that choice depends on
continuous financial quantities like covariance and expected return.

\subsection{Why Quantum?}

Combinatorial optimization over binary variables is a natural fit for quantum
computing. The problem can be encoded as a QUBO (Quadratic Unconstrained Binary
Optimization), which maps directly to an Ising Hamiltonian, which is exactly
what quantum algorithms like QAOA (Quantum Approximate Optimization Algorithm)
are designed to minimize.

At 8 qubits our problem is small enough that a classical computer can solve it
exactly in microseconds. The point of this project is not to demonstrate quantum
advantage. It is to:

\begin{enumerate}
    \item Build a correct, end-to-end quantum optimization pipeline on real insurance data.
    \item Analyze how quantum hardware properties (connectivity, noise) affect solution quality.
    \item Demonstrate the workflow that \textit{would} provide value at scale (16--50+ qubits).
\end{enumerate}

\subsection{The Data}

The Hartford provides an Excel workbook containing:

\begin{itemize}
    \item \textbf{Assets sheet}: 50 assets, each with a sector label, expected return, capital charge, and liquidity score.
    \item \textbf{Scenarios sheet}: 1{,}200 economic scenarios, each giving a return value for every asset.
    \item \textbf{Covariance sheet}: A $50 \times 50$ asset-level covariance matrix.
\end{itemize}

The 50 assets belong to 8 natural sectors:
\begin{center}
\begin{tabular}{llll}
\toprule
Cash & Gov Bonds & IG Credit & HY Credit \\
Equities US & Equities Intl & Real Estate & Infrastructure \\
\bottomrule
\end{tabular}
\end{center}

\subsection{The Decision}

We select exactly $B = 4$ sectors out of 8. Each selected sector receives an
equal weight of $1/B = 0.25$. This gives us a binary decision vector
$\mathbf{x} \in \{0,1\}^8$ with the constraint $\sum_i x_i = B$.

The number of feasible portfolios is:
\[
\binom{8}{4} = 70
\]

The total search space (including infeasible) is $2^8 = 256$.


% ══════════════════════════════════════════════════════════════════════════════
\section{Background Theory: Portfolio Optimization}
% ══════════════════════════════════════════════════════════════════════════════

\subsection{Mean--Variance Optimization (Markowitz)}

Modern portfolio theory, introduced by Harry Markowitz in 1952, frames
portfolio construction as a trade-off between expected return and risk
(variance).

\begin{definition}[Portfolio Return and Risk]
Given a weight vector $\mathbf{w} \in \mathbb{R}^n$ over $n$ assets with
expected return vector $\bm{\mu} \in \mathbb{R}^n$ and covariance matrix
$\bm{\Sigma} \in \mathbb{R}^{n \times n}$:

\begin{align}
\text{Expected return:} \quad & R_p = \bm{\mu}^\top \mathbf{w} \\[4pt]
\text{Portfolio variance:} \quad & \sigma_p^2 = \mathbf{w}^\top \bm{\Sigma}\, \mathbf{w} \\[4pt]
\text{Portfolio volatility:} \quad & \sigma_p = \sqrt{\mathbf{w}^\top \bm{\Sigma}\, \mathbf{w}}
\end{align}
\end{definition}

The classical Markowitz problem is:
\[
\min_{\mathbf{w}} \; q \cdot \mathbf{w}^\top \bm{\Sigma}\, \mathbf{w}
    - \bm{\mu}^\top \mathbf{w}
\quad \text{subject to} \quad
\sum_i w_i = 1, \;\; w_i \geq 0
\]

where $q > 0$ is the \textbf{risk aversion parameter}. A higher $q$ penalizes
variance more heavily, producing more conservative portfolios.

\begin{remark}
In our project, the weights are not free continuous variables. Because we use
equal weighting among selected sectors, the weight vector is fully determined
by the binary selection vector: $w_i = x_i / B$.
\end{remark}

\subsection{Why Covariance Matters}

The covariance matrix $\bm{\Sigma}$ captures how asset returns move together.
Its role in portfolio optimization is crucial:

\begin{itemize}
    \item \textbf{Diagonal entries} $\Sigma_{ii}$ are individual sector variances.
    \item \textbf{Off-diagonal entries} $\Sigma_{ij}$ measure co-movement.
    \item Selecting sectors with \textit{low} or \textit{negative} covariance produces diversification, reducing overall portfolio risk.
\end{itemize}

The quadratic term $\mathbf{x}^\top \bm{\Sigma}\, \mathbf{x}$ in our objective
is what makes this problem naturally expressible as a QUBO: it involves pairwise
interactions between binary variables, exactly the structure a QUBO encodes.

\subsection{Insurance-Specific Metrics}

Beyond return and variance, insurance companies care about:

\begin{description}[leftmargin=2cm]
    \item[Capital Charge] Regulatory frameworks (e.g., Solvency~II) require insurers
    to hold capital reserves proportional to investment risk. A high capital charge
    means the insurer must set aside more capital, reducing the amount available
    for other uses.

    \item[Liquidity] Insurers must be able to pay claims. Illiquid investments
    (e.g., real estate, infrastructure) may offer higher returns but cannot be
    sold quickly in a crisis.

    \item[VaR (Value at Risk)] The loss level that is exceeded with probability
    $1 - \alpha$. At $\alpha = 0.95$:
    \[
    \text{VaR}_{0.95} = \inf\{l : P(\text{Loss} > l) \leq 0.05\}
    \]

    \item[CVaR (Conditional Value at Risk)] The expected loss given that the loss
    exceeds VaR. Also called Expected Shortfall:
    \[
    \text{CVaR}_\alpha = \mathbb{E}[\text{Loss} \mid \text{Loss} \geq \text{VaR}_\alpha]
    \]
    CVaR is generally considered a better risk measure than VaR because it
    accounts for the severity of tail losses, not just their frequency.
\end{description}


% ══════════════════════════════════════════════════════════════════════════════
\section{Data Pipeline: From Workbook to 8 Sectors}
% ══════════════════════════════════════════════════════════════════════════════

\subsection{Why Aggregate?}

The raw workbook has 50 assets but we want 8 qubits (one per sector). We need
to reduce the dimensionality while preserving as much of the real financial
structure as possible.

\subsection{The Aggregation Matrix}

Let $R \in \mathbb{R}^{S \times N}$ be the scenario-return matrix ($S = 1200$
scenarios, $N = 50$ assets). We construct an \textbf{aggregation matrix}
$A \in \mathbb{R}^{8 \times 50}$ where:

\[
A_{s,a} = \begin{cases}
\dfrac{1}{|\mathcal{A}_s|} & \text{if asset } a \text{ belongs to sector } s \\[6pt]
0 & \text{otherwise}
\end{cases}
\]

where $\mathcal{A}_s$ is the set of assets in sector $s$. This gives each
asset within a sector an equal weight, producing an average sector return.

The sector-level scenario returns are:
\[
R_{\text{sec}} = R \cdot A^\top \in \mathbb{R}^{S \times 8}
\]

From this matrix we compute:
\begin{align}
\bm{\mu} &= \frac{1}{S} \sum_{s=1}^{S} (R_{\text{sec}})_s \in \mathbb{R}^8
    \quad \text{(sector expected returns)} \\[6pt]
\bm{\Sigma} &= \text{Cov}(R_{\text{sec}}) \in \mathbb{R}^{8 \times 8}
    \quad \text{(sector covariance matrix)}
\end{align}

\subsection{Capital and Liquidity Vectors}

For each sector, we average the capital charges and liquidity scores across all
assets in that sector:

\[
c_s = \frac{1}{|\mathcal{A}_s|} \sum_{a \in \mathcal{A}_s} \text{capital\_charge}(a)
\qquad
\ell_s = \frac{1}{|\mathcal{A}_s|} \sum_{a \in \mathcal{A}_s} \text{liquidity}(a)
\]

These are then normalized to $[0, 1]$ using min--max normalization:
\[
\tilde{c}_s = \frac{c_s - \min(\mathbf{c})}{\max(\mathbf{c}) - \min(\mathbf{c})}
\qquad
\tilde{\imath}_s = 1 - \frac{\ell_s - \min(\bm{\ell})}{\max(\bm{\ell}) - \min(\bm{\ell})}
\]

Note that illiquidity is $1 - \text{normalized liquidity}$, because we want to
\textit{penalize} illiquid sectors (the higher the illiquidity, the worse).

\subsection{Validation}

We cross-check our aggregated quantities against the workbook's own covariance
sheet by computing:
\[
\text{Frobenius error} = \| \bm{\Sigma}_{\text{scenarios}} - A\, \bm{\Sigma}_{\text{workbook}}\, A^\top \|_F
\]

In our run this was $\approx 0.040$, confirming that the scenario-based and
workbook-based covariance estimates are consistent.


% ══════════════════════════════════════════════════════════════════════════════
\section{QUBO Formulation}
% ══════════════════════════════════════════════════════════════════════════════

\subsection{What is a QUBO?}

\begin{definition}[QUBO]
A Quadratic Unconstrained Binary Optimization problem has the form:
\[
\min_{\mathbf{x} \in \{0,1\}^n} \; \mathbf{x}^\top Q\, \mathbf{x}
\]
where $Q \in \mathbb{R}^{n \times n}$ is an upper-triangular matrix encoding
both linear terms (on the diagonal) and quadratic interactions (off-diagonal).
\end{definition}

The key property is that for binary variables, $x_i^2 = x_i$, so the diagonal
of $Q$ captures linear costs and the off-diagonal captures pairwise interactions.

\subsection{Our Objective Function}

We build four model variants. The full insurer-aware objective (``both'' model) is:

\begin{equation}
\boxed{
\min_{\mathbf{x} \in \{0,1\}^8} \;
\underbrace{\frac{q}{B^2}\, \mathbf{x}^\top \bm{\Sigma}\, \mathbf{x}}_{\text{risk}}
\;-\; \underbrace{\frac{1}{B}\, \bm{\mu}^\top \mathbf{x}}_{\text{return}}
\;+\; \underbrace{\lambda\, \bigl(\mathbf{1}^\top \mathbf{x} - B\bigr)^2}_{\text{budget penalty}}
\;+\; \underbrace{\frac{\gamma}{B}\, \tilde{\mathbf{c}}^\top \mathbf{x}}_{\text{capital}}
\;+\; \underbrace{\frac{\eta}{B}\, \tilde{\bm{\imath}}^\top \mathbf{x}}_{\text{illiquidity}}
}
\label{eq:full_objective}
\end{equation}

where:

\begin{center}
\begin{tabular}{clc}
\toprule
\textbf{Symbol} & \textbf{Meaning} & \textbf{Default Value} \\
\midrule
$\mathbf{x}$ & Binary sector selection vector & --- \\
$B$ & Number of sectors to select & 4 \\
$q$ & Risk aversion & 0.5 \\
$\lambda$ & Budget penalty strength & 8.0 \\
$\gamma$ & Capital charge penalty & 0.25 \\
$\eta$ & Illiquidity penalty & 0.25 \\
$\bm{\mu}$ & Sector expected return vector & from data \\
$\bm{\Sigma}$ & Sector covariance matrix & from data \\
$\tilde{\mathbf{c}}$ & Normalized capital charges & from data \\
$\tilde{\bm{\imath}}$ & Normalized illiquidity scores & from data \\
\bottomrule
\end{tabular}
\end{center}

\subsection{The Four Model Variants}

\begin{center}
\begin{tabular}{lccl}
\toprule
\textbf{Variant} & $\gamma$ & $\eta$ & \textbf{What It Tests} \\
\midrule
\texttt{base} & 0 & 0 & Pure Markowitz risk--return \\
\texttt{capital} & 0.25 & 0 & Adds regulatory capital cost \\
\texttt{liquidity} & 0 & 0.25 & Adds illiquidity cost \\
\texttt{both} & 0.25 & 0.25 & Full insurer-aware model \\
\bottomrule
\end{tabular}
\end{center}

Running all four lets us show judges exactly how each insurance constraint
changes the optimal portfolio.

\subsection{Expanding the Budget Penalty}

The budget constraint $\sum_i x_i = B$ is enforced via a quadratic penalty:
\[
\lambda\, \bigl(\mathbf{1}^\top \mathbf{x} - B\bigr)^2
= \lambda \left( \sum_i x_i - B \right)^2
\]

Expanding this:
\begin{align}
\lambda \left(\sum_i x_i - B\right)^2
&= \lambda \left( \sum_i x_i^2 + 2\sum_{i<j} x_i x_j - 2B\sum_i x_i + B^2 \right) \\
&= \lambda \left( \sum_i x_i + 2\sum_{i<j} x_i x_j - 2B\sum_i x_i + B^2 \right)
    \quad \text{(since } x_i^2 = x_i \text{)} \\
&= \lambda \sum_i (1 - 2B)\, x_i + 2\lambda \sum_{i<j} x_i x_j + \lambda B^2
\end{align}

The constant $\lambda B^2$ does not affect the optimization and can be dropped.

\subsection{Constructing the QUBO Matrix}

Combining all terms, the upper-triangular QUBO matrix $Q$ has entries:

\textbf{Diagonal entries} ($i = j$):
\begin{equation}
Q_{ii} = \frac{q}{B^2}\, \Sigma_{ii}
         - \frac{1}{B}\, \mu_i
         + \lambda\,(1 - 2B)
         + \frac{\gamma}{B}\, \tilde{c}_i
         + \frac{\eta}{B}\, \tilde{\imath}_i
\label{eq:qubo_diagonal}
\end{equation}

\textbf{Off-diagonal entries} ($i < j$):
\begin{equation}
Q_{ij} = \frac{2q}{B^2}\, \Sigma_{ij} + 2\lambda
\label{eq:qubo_offdiag}
\end{equation}

\begin{remark}
The factor of 2 on both $\Sigma_{ij}$ and $\lambda$ comes from the fact that
we store $Q$ as upper-triangular: the contribution of the pair $(i,j)$ appears
only once in $Q_{ij}$ rather than being split between $Q_{ij}$ and $Q_{ji}$.
\end{remark}

\subsection{QUBO Scaling}

The entries of $Q$ can span a wide numerical range. Before converting to an
Ising Hamiltonian, we normalize:
\[
Q_{\text{scaled}} = \frac{Q}{\alpha}, \qquad \alpha = \max_{i,j} |Q_{ij}|
\]

This brings all matrix entries into $[-1, 1]$, which improves the numerical
conditioning of the QAOA simulation. The unscaled $Q$ is kept for computing
physically meaningful portfolio objectives.

\subsection{Why $\lambda = 8.0$?}

The penalty $\lambda$ must be large enough that any infeasible bitstring
(one that does not select exactly $B$ sectors) has a higher objective than
the best feasible bitstring. If $\lambda$ is too small, the QUBO global
minimum might be an infeasible solution.

Our code automatically verifies this: after constructing the QUBO, it checks
whether the brute-force global optimum is feasible. If not, it doubles $\lambda$
and re-checks, repeating up to 6 times. With $\lambda = 8.0$, all four model
variants have feasible global optima without needing adjustment.

\begin{remark}
The penalty approach converts a \textit{constrained} problem into an
\textit{unconstrained} one. This is necessary because QUBO and Ising
formulations do not natively support hard constraints. The trade-off is that
very large $\lambda$ values can distort the energy landscape, making the
penalty-related quadratic terms dominate over the actual portfolio quality terms.
\end{remark}


% ══════════════════════════════════════════════════════════════════════════════
\section{Ising Hamiltonian}
% ══════════════════════════════════════════════════════════════════════════════

\subsection{Why Convert to Ising?}

Quantum algorithms like QAOA work with quantum mechanical Hamiltonians, not
QUBO matrices. The standard bridge is the Ising model:

\begin{definition}[Ising Hamiltonian]
The classical Ising Hamiltonian for $n$ spins is:
\[
H_C = \sum_{i=1}^{n} h_i\, Z_i + \sum_{i < j} J_{ij}\, Z_i\, Z_j
\]
where $Z_i \in \{+1, -1\}$ are spin variables (eigenvalues of the Pauli-$Z$ operator),
$h_i$ are local fields, and $J_{ij}$ are coupling strengths.
\end{definition}

\subsection{The Mapping: Binary to Spin}

The standard substitution is:
\begin{equation}
x_i = \frac{1 - s_i}{2}
\label{eq:bit_to_spin}
\end{equation}

where $s_i \in \{+1, -1\}$ and $x_i \in \{0, 1\}$. This gives:
\begin{align}
x_i = 0 &\;\longleftrightarrow\; s_i = +1 \\
x_i = 1 &\;\longleftrightarrow\; s_i = -1
\end{align}

\begin{remark}
This is the \textit{opposite} of the common physics convention where
$s_i = +1$ means ``spin up'' (selected). Our convention comes from the
standard QUBO-to-Ising literature and is used consistently throughout the
code. If a judge asks, we can explain this: it does not affect correctness
as long as the mapping is applied uniformly.
\end{remark}

\subsection{Deriving the Ising Coefficients}

Substituting $x_i = (1 - s_i)/2$ into $\mathbf{x}^\top Q\,\mathbf{x}$:

\textbf{Diagonal terms} ($Q_{ii} x_i^2 = Q_{ii} x_i$):
\[
Q_{ii}\, x_i = Q_{ii}\, \frac{1 - s_i}{2} = \frac{Q_{ii}}{2} - \frac{Q_{ii}}{2}\, s_i
\]

This contributes $Q_{ii}/2$ to the constant offset and $-Q_{ii}/2$ to $h_i$.

\textbf{Off-diagonal terms} ($Q_{ij} x_i x_j$ for $i < j$):
\begin{align}
Q_{ij}\, x_i\, x_j
&= Q_{ij}\, \frac{1 - s_i}{2}\, \frac{1 - s_j}{2} \\
&= \frac{Q_{ij}}{4}\, (1 - s_i - s_j + s_i s_j)
\end{align}

This contributes:
\begin{itemize}
    \item $Q_{ij}/4$ to the constant offset
    \item $-Q_{ij}/4$ to $h_i$
    \item $-Q_{ij}/4$ to $h_j$
    \item $+Q_{ij}/4$ to $J_{ij}$
\end{itemize}

\textbf{Summary of the Ising coefficients:}
\begin{align}
h_i &= -\frac{Q_{ii}}{2} - \sum_{j \neq i} \frac{Q_{\min(i,j),\max(i,j)}}{4}
\label{eq:ising_h} \\[6pt]
J_{ij} &= \frac{Q_{ij}}{4} \quad \text{for } i < j
\label{eq:ising_J} \\[6pt]
\text{offset} &= \sum_i \frac{Q_{ii}}{2} + \sum_{i<j} \frac{Q_{ij}}{4}
\label{eq:ising_offset}
\end{align}

The offset is a constant that shifts the entire energy spectrum. It does not
affect which bitstring is optimal, but we track it for accurate energy
comparisons.


% ══════════════════════════════════════════════════════════════════════════════
\section{Classical Baselines}
% ══════════════════════════════════════════════════════════════════════════════

We use three classical methods to validate our quantum results. These are not
competitors---they are ground truth.

\subsection{Exact Brute-Force Enumeration}

With $n = 8$ qubits, the total search space is $2^8 = 256$ bitstrings. We
evaluate the QUBO objective for every single one:

\[
\text{For each } \mathbf{x} \in \{0,1\}^8: \quad
f(\mathbf{x}) = \mathbf{x}^\top Q\, \mathbf{x}
\]

We then find:
\begin{itemize}
    \item The \textbf{best overall} bitstring (lowest $f$, may be infeasible).
    \item The \textbf{best feasible} bitstring (lowest $f$ among those with $\sum x_i = B$).
\end{itemize}

If the best overall is also feasible, our $\lambda$ is large enough. If not,
we double $\lambda$ and repeat.

This is computationally trivial at $n = 8$ but would become impossible at
$n \geq 30$ ($2^{30} \approx 10^9$ evaluations). That is where quantum
methods would become valuable.

\subsection{Greedy Selection}

The greedy baseline adds one sector at a time, always picking the sector that
most reduces the QUBO objective:

\begin{algorithm}[H]
\caption{Greedy Sector Selection}
\begin{algorithmic}[1]
\State $\mathbf{x} \gets \mathbf{0}$, $\text{selected} \gets \emptyset$
\For{$k = 1$ to $B$}
    \State $i^* \gets \arg\min_{i \notin \text{selected}} f(\mathbf{x} + \mathbf{e}_i)$
    \State $x_{i^*} \gets 1$, $\text{selected} \gets \text{selected} \cup \{i^*\}$
\EndFor
\State \Return $\mathbf{x}$
\end{algorithmic}
\end{algorithm}

In all four model variants, greedy finds the exact optimum. This is not
guaranteed in general---greedy can get trapped in local optima---but it happens
to work here because the 8-sector problem has relatively smooth structure.

\subsection{Continuous Markowitz Baseline}

We also solve the \textit{continuous relaxation}: instead of binary $x_i$,
we allow real-valued weights $w_i \geq 0$ with $\sum w_i = 1$.

The continuous problem is:
\[
\min_{\mathbf{w} \in \Delta} \; q\, \mathbf{w}^\top \bm{\Sigma}\, \mathbf{w}
    - \bm{\mu}^\top \mathbf{w}
\]

where $\Delta = \{\mathbf{w} : w_i \geq 0, \sum w_i = 1\}$ is the probability simplex.

We solve this with projected gradient descent:

\begin{algorithm}[H]
\caption{Continuous Markowitz via Projected Gradient Descent}
\begin{algorithmic}[1]
\State $\mathbf{w} \gets (1/8, \ldots, 1/8)$
\For{$t = 1$ to $T$}
    \State $\mathbf{g} \gets 2q\, \bm{\Sigma}\, \mathbf{w} - \bm{\mu}$
    \State $\eta_t \gets \eta_0 / \sqrt{t}$ \Comment{Decaying learning rate}
    \State $\mathbf{w} \gets \text{proj}_\Delta(\mathbf{w} - \eta_t\, \mathbf{g})$
\EndFor
\end{algorithmic}
\end{algorithm}

The simplex projection uses the algorithm of Duchi et al.\ (2008): sort the
components, find a threshold, and clip.

After convergence, we round to a discrete portfolio by selecting the top-$B$
sectors by weight. This ``continuous proxy'' is then evaluated with the full
QUBO objective for comparison.


% ══════════════════════════════════════════════════════════════════════════════
\section{QAOA: Quantum Approximate Optimization Algorithm}
% ══════════════════════════════════════════════════════════════════════════════

\subsection{The Idea}

QAOA, introduced by Farhi, Goldstone, and Gutmann (2014), is a variational
quantum algorithm designed for combinatorial optimization. It constructs a
parameterized quantum state whose measurement outcomes are biased toward
low-energy (good) solutions of the cost Hamiltonian.

\subsection{The QAOA Circuit}

For depth $p$, QAOA alternates between two operations:

\begin{enumerate}
    \item \textbf{Cost layer}: Apply $e^{-i \gamma_k H_C}$ where $H_C$ is the Ising cost Hamiltonian.
    \item \textbf{Mixer layer}: Apply $e^{-i \beta_k H_M}$ where $H_M = \sum_i X_i$ is the transverse-field mixer.
\end{enumerate}

The full QAOA state is:
\begin{equation}
\boxed{
|\bm{\gamma}, \bm{\beta}\rangle
= \prod_{k=1}^{p} e^{-i \beta_k H_M}\, e^{-i \gamma_k H_C}\, |+\rangle^{\otimes n}
}
\label{eq:qaoa_state}
\end{equation}

where $|+\rangle = \frac{1}{\sqrt{2}}(|0\rangle + |1\rangle)$ and the initial
state $|+\rangle^{\otimes n}$ is a uniform superposition over all $2^n$
computational basis states.

\subsection{The Cost Layer in Detail}

Since $H_C$ is diagonal in the computational basis, $e^{-i \gamma H_C}$
applies a phase to each basis state proportional to its energy:

\[
e^{-i \gamma H_C} |z\rangle = e^{-i \gamma E(z)} |z\rangle
\]

where $E(z)$ is the Ising energy of bitstring $z$. In our implementation, we
precompute all 256 energies and apply the phases as element-wise multiplication
on the state vector.

\subsection{The Mixer Layer in Detail}

The mixer Hamiltonian $H_M = \sum_i X_i$ generates rotations around the
$X$-axis on each qubit independently:

\[
e^{-i \beta H_M} = \bigotimes_{i=1}^{n} R_X(2\beta)
\]

where $R_X(\theta) = e^{-i \theta X / 2}$. In matrix form:

\[
R_X(\theta) = \begin{pmatrix}
\cos(\theta/2) & -i\sin(\theta/2) \\
-i\sin(\theta/2) & \cos(\theta/2)
\end{pmatrix}
\]

Our code applies this qubit-by-qubit on the state vector using stride-based
indexing, without ever building the full $256 \times 256$ matrix.

\subsection{Variational Optimization}

The $2p$ parameters $(\gamma_1, \ldots, \gamma_p, \beta_1, \ldots, \beta_p)$
are optimized classically to minimize the expected cost:

\[
C(\bm{\gamma}, \bm{\beta})
= \langle \bm{\gamma}, \bm{\beta} | H_C | \bm{\gamma}, \bm{\beta} \rangle
= \sum_{z} |\langle z | \bm{\gamma}, \bm{\beta} \rangle|^2 \cdot E(z)
\]

Our optimizer uses \textbf{multi-start coordinate descent}:

\begin{enumerate}
    \item \textbf{Global exploration}: Sample 12 random starting points.
    For each, evaluate the expected energy. Keep the best.
    \item \textbf{Local refinement}: Starting from the best random point,
    perturb each parameter in both directions by a step size $\delta$. If any
    perturbation improves the energy, accept it. After each full pass, halve
    the step size. Repeat for up to 5 rounds or until no improvement is found.
\end{enumerate}

\begin{remark}
This optimizer is simple and dependency-free, which was a deliberate choice
for hackathon robustness. A stronger optimizer (e.g., COBYLA or Nelder--Mead
from SciPy) would likely find slightly better angles, but at 8 qubits the
improvement would be marginal and would not change our conclusions.
\end{remark}

\subsection{Why QAOA Depths $p = 1, 2, 3$?}

\begin{itemize}
    \item $p = 1$: The shallowest meaningful QAOA. Has $2$ parameters. Produces
    a smooth energy landscape that is easy to visualize and optimize, but has
    limited expressivity.
    \item $p = 2$: Four parameters. Can capture more complex correlations
    between the cost and mixer layers.
    \item $p = 3$: Six parameters. Deeper circuits can in principle better
    approximate the ground state, but optimization becomes harder and hardware
    noise accumulates.
\end{itemize}

In the limit $p \to \infty$, QAOA can find the exact ground state (by the
quantum adiabatic theorem), but practical quantum hardware limits $p$ to
single digits.

\subsection{Measurement and Sampling}

After preparing the QAOA state, we obtain a probability distribution over all
$2^n$ basis states. To simulate a real quantum device, we draw $2{,}048$
samples (``shots'') from this distribution and count occurrences of each
bitstring.

Key metrics from the output distribution:

\begin{align}
P(\text{optimal}) &= |\langle z^* | \bm{\gamma}, \bm{\beta}\rangle|^2 \\
P(\text{feasible}) &= \sum_{z : |z| = B} |\langle z | \bm{\gamma}, \bm{\beta}\rangle|^2 \\
P(\text{optimal} \mid \text{feasible}) &= \frac{P(\text{optimal})}{P(\text{feasible})}
\end{align}

where $z^*$ is the known classical optimum and $|z| = \sum_i z_i$.


% ══════════════════════════════════════════════════════════════════════════════
\section{Circuit Construction for Neutral-Atom Hardware}
% ══════════════════════════════════════════════════════════════════════════════

\subsection{From Ising to Gates}

The cost layer requires implementing $e^{-i \gamma H_C}$. Since:
\[
H_C = \sum_i h_i Z_i + \sum_{i<j} J_{ij} Z_i Z_j
\]

the exponential factorizes (all terms commute because they are all diagonal):

\[
e^{-i \gamma H_C} = \prod_i e^{-i \gamma h_i Z_i}
                     \prod_{i<j} e^{-i \gamma J_{ij} Z_i Z_j}
\]

\textbf{Single-qubit terms}: $e^{-i \gamma h_i Z_i} = R_Z(2\gamma h_i)$.

\textbf{Two-qubit terms}: $e^{-i \gamma J_{ij} Z_i Z_j}$ requires the
decomposition:

\begin{equation}
e^{-i \theta Z_i Z_j} = \text{CNOT}(i,j) \;\cdot\; R_Z(2\theta)_j \;\cdot\; \text{CNOT}(i,j)
\label{eq:zz_decomp}
\end{equation}

\begin{remark}
An earlier version of the code used an H/CZ/RZ/CZ/H decomposition, which does
\textbf{not} implement $e^{-i\theta ZZ}$ correctly. This was a real bug that
caused a mismatch between the NumPy simulator and the Cirq circuit path. The
fix was to switch to the CNOT/RZ/CNOT decomposition shown above, which is the
standard correct form. This is documented in the code comments.
\end{remark}

\subsection{Gate Scheduling and Connectivity}

On a real neutral-atom device (like QuEra's Bloqade), not all qubit pairs can
interact simultaneously. The circuit depth depends on how we schedule the
two-qubit gates.

We test four \textbf{connectivity modes}:

\begin{description}[leftmargin=2.5cm]
    \item[naive] Each ZZ interaction is applied sequentially, one at a time.
    For $K_8$ (complete graph on 8 vertices, 28 edges), this gives 28 sequential
    groups, each containing 3 gate layers (CNOT--RZ--CNOT), yielding the deepest
    circuit.

    \item[hand] Use a greedy edge coloring algorithm to partition the 28 edges
    into disjoint matchings. Each matching can be executed in parallel. For $K_8$,
    the chromatic index is 7, so we get 7 groups instead of 28. This drastically
    reduces circuit depth.

    \item[auto] Use Bloqade's built-in parallelization pass (when installed) to
    automatically optimize the circuit layout. Falls back to \texttt{hand} if
    Bloqade is not available.

    \item[pruned] Only keep the 12 strongest (by $|J_{ij}|$) coupling edges. This
    reduces the two-qubit gate count at the cost of losing weaker interactions.
\end{description}

\subsection{Edge Coloring}

\begin{definition}[Edge Coloring]
An edge coloring of a graph $G$ assigns colors to edges such that no two edges
sharing a vertex have the same color. The minimum number of colors needed is
the \textbf{chromatic index} $\chi'(G)$.
\end{definition}

For a complete graph $K_n$ with $n$ even: $\chi'(K_n) = n - 1$. So $K_8$ needs
7 colors, meaning 7 parallel groups of 4 simultaneous ZZ gates each.

Our implementation uses a greedy algorithm: for each edge, try to assign it to
an existing color class (checking that no vertex conflicts exist). If no class
works, create a new one.

\subsection{Circuit Depth Comparison}

\begin{center}
\begin{tabular}{lccc}
\toprule
\textbf{Mode} & \textbf{ZZ Groups} & \textbf{Gates per layer} & \textbf{Relative Depth} \\
\midrule
\texttt{naive} & 28 & $28 \times 3 = 84$ & Deepest \\
\texttt{hand} / \texttt{auto} & 7 & $7 \times 3 = 21$ & $4\times$ reduction \\
\texttt{pruned} & $\leq 7$ & $\leq 21$ & Shallowest \\
\bottomrule
\end{tabular}
\end{center}

\subsection{The Atom Layout Figure}

To give a neutral-atom-specific perspective, we compute a 2D layout of the 8
sectors using classical Multidimensional Scaling (MDS) on the correlation
distance matrix:

\[
D_{ij} = 1 - |\rho_{ij}|
\]

where $\rho_{ij} = \Sigma_{ij} / (\sigma_i \sigma_j)$ is the Pearson
correlation between sectors $i$ and $j$. Sectors with high correlation are
placed close together, mimicking how atoms would be arranged to minimize
long-range interactions.


% ══════════════════════════════════════════════════════════════════════════════
\section{Noise Models}
% ══════════════════════════════════════════════════════════════════════════════

\subsection{Why Study Noise?}

Real quantum devices are noisy. Gates have finite fidelity, qubits decohere,
and measurement is imperfect. Understanding how noise degrades our results is
essential for assessing whether this workflow would work on real hardware.

\subsection{Our Approximate Noise Model}

When Bloqade's hardware-calibrated noise transforms are not available, we use
a simplified model that mixes the ideal probability distribution with the
uniform distribution:

\begin{equation}
\mathbf{p}_{\text{noisy}} = (1 - \epsilon)\, \mathbf{p}_{\text{ideal}}
    + \epsilon\, \mathbf{u}
\label{eq:noise_mixing}
\end{equation}

where $\mathbf{u} = (1/2^n, \ldots, 1/2^n)$ is the uniform distribution and
the effective noise strength is:

\[
\epsilon = \min\bigl(0.95,\; \alpha_{\text{model}} \cdot s
    \cdot (0.0015 \cdot d + 0.002 \cdot g)\bigr)
\]

with $d$ = circuit depth, $g$ = two-qubit gate count, $s$ = user-specified noise
strength, and $\alpha_{\text{model}}$ a per-model scale factor.

\subsection{Noise Model Variants}

\begin{center}
\begin{tabular}{lcp{7cm}}
\toprule
\textbf{Model} & $\alpha$ & \textbf{Physical Motivation} \\
\midrule
\texttt{none} & 0.0 & Ideal noiseless simulation \\
\texttt{one\_zone} & 1.0 & Single atom-loss zone (dominant neutral-atom error) \\
\texttt{two\_zone} & 0.75 & Two independent loss zones (partially correlated) \\
\texttt{depolarizing} & 1.25 & Single-qubit depolarizing channel (worst case) \\
\bottomrule
\end{tabular}
\end{center}

\begin{remark}
We are transparent that these noise models are \textbf{not} hardware-calibrated.
They all use the same mixing formula with different scale factors. The purpose
is to demonstrate the \textit{type} of analysis (how does solution quality
degrade with noise?) rather than to predict exact hardware performance. When
Bloqade is installed, the code automatically upgrades to use Bloqade's native
\texttt{GeminiOneZoneNoiseModel} and \texttt{GeminiTwoZoneNoiseModel}.
\end{remark}

\subsection{Readout Error}

Independent of the noise model, we also simulate measurement errors. Each qubit
has a probability $p_{\text{ro}}$ of being flipped during readout:

\[
P(\text{measure } z' \mid \text{true state } z) = \prod_{i=1}^n
\begin{cases}
1 - p_{\text{ro}} & \text{if } z'_i = z_i \\
p_{\text{ro}} & \text{if } z'_i \neq z_i
\end{cases}
\]

This is applied to the full probability vector:
\[
p_{\text{measured}}(z') = \sum_{z} p_{\text{true}}(z) \cdot (1 - p_{\text{ro}})^{n - d_H(z, z')} \cdot p_{\text{ro}}^{d_H(z, z')}
\]

where $d_H(z, z')$ is the Hamming distance between $z$ and $z'$.


% ══════════════════════════════════════════════════════════════════════════════
\section{Results}
% ══════════════════════════════════════════════════════════════════════════════

\subsection{Classical Optima}

The exact brute-force search on the real Hartford workbook produces:

\begin{center}
\begin{tabular}{llp{6cm}}
\toprule
\textbf{Model} & \textbf{Bitstring} & \textbf{Selected Sectors} \\
\midrule
\texttt{base} & \texttt{00011110} & HY Credit, Equities US, Equities Intl, Real Estate \\
\texttt{capital} & \texttt{11110000} & Cash, Gov Bonds, IG Credit, HY Credit \\
\texttt{liquidity} & \texttt{11001100} & Cash, Gov Bonds, Equities US, Equities Intl \\
\texttt{both} & \texttt{11110000} & Cash, Gov Bonds, IG Credit, HY Credit \\
\bottomrule
\end{tabular}
\end{center}

\textbf{Business interpretation:}
\begin{itemize}
    \item \texttt{base} selects high-return sectors (equities, HY credit, real estate), accepting higher risk and capital cost.
    \item \texttt{capital} rotates entirely into low-capital-charge fixed income, sacrificing return for regulatory efficiency.
    \item \texttt{liquidity} keeps some equity exposure but drops illiquid sectors (real estate, infrastructure).
    \item \texttt{both} matches \texttt{capital}, indicating that at these penalty levels ($\gamma = \eta = 0.25$), the capital constraint dominates the liquidity constraint.
\end{itemize}

\subsection{Quantum Results}

\begin{center}
\begin{tabular}{lcccc}
\toprule
\textbf{Metric} & \textbf{Base} & \textbf{Capital} & \textbf{Liquidity} & \textbf{Both} \\
\midrule
Quantum bitstring matches exact? & Yes & Yes & Yes & Yes \\
$P(\text{optimal})$ & 1.03\% & 1.03\% & 1.03\% & 1.03\% \\
$P(\text{feasible})$ & 71.96\% & 71.96\% & 71.96\% & 71.96\% \\
$P(\text{optimal} \mid \text{feasible})$ & 1.43\% & 1.43\% & 1.43\% & 1.43\% \\
Uniform baseline $1/\binom{8}{4}$ & 1.43\% & 1.43\% & 1.43\% & 1.43\% \\
\bottomrule
\end{tabular}
\end{center}

\subsection{Interpreting the Quantum Performance}

The probability $P(\text{optimal} \mid \text{feasible}) \approx 1/70$ means
QAOA is essentially \textbf{uniform over the feasible set}. It learns the budget
constraint ($\sum x_i = 4$) but does not distinguish between the 70 feasible
portfolios.

\textbf{Why?} The energy landscape figure shows the answer: the gap between the
best and second-best feasible bitstring is only $0.000031$ in the scaled QUBO.
With such a tiny gap, even an ideal QAOA at shallow depth cannot concentrate
significant probability mass on the optimum.

This is \textit{not} a failure of our implementation. It is a fundamental
property of this particular problem instance at this scale. The practical
implication is that at 8 qubits, brute-force enumeration is both faster and
more reliable.

\subsection{Connectivity Results}

\begin{center}
\begin{tabular}{lcc}
\toprule
\textbf{Mode} & \textbf{Avg.\ Circuit Depth} & \textbf{Avg.\ $P(\text{opt})$} \\
\midrule
\texttt{naive} & $\sim170$ & $\sim0.009$ \\
\texttt{auto} & $\sim60$ & $\sim0.009$ \\
\texttt{hand} & $\sim60$ & $\sim0.009$ \\
\texttt{pruned} & $\sim45$ & $\sim0.004$ \\
\bottomrule
\end{tabular}
\end{center}

Key observations:
\begin{itemize}
    \item \texttt{hand} and \texttt{auto} achieve the same quality as \texttt{naive} at $\sim3\times$ less depth. This is pure scheduling gain.
    \item \texttt{pruned} halves the gate count but loses solution quality because the discarded weak couplings still carry relevant information.
    \item At noiseless simulation, depth does not hurt (no decoherence). On real hardware, the depth reduction would translate directly to improved fidelity.
\end{itemize}

\subsection{Noise Results}

Under increasing noise strength:
\begin{itemize}
    \item All noise models degrade $P(\text{optimal})$ monotonically.
    \item \texttt{depolarizing} ($\alpha = 1.25$) degrades fastest.
    \item \texttt{two\_zone} ($\alpha = 0.75$) degrades slowest.
    \item At noise strength 2.0, the optimal probability drops by $\sim30$--$40\%$ from the noiseless value.
\end{itemize}


% ══════════════════════════════════════════════════════════════════════════════
\section{Code Architecture}
% ══════════════════════════════════════════════════════════════════════════════

\subsection{File Overview}

\begin{center}
\begin{longtable}{lcp{7cm}}
\toprule
\textbf{File} & \textbf{Lines} & \textbf{Purpose} \\
\midrule
\endhead
\texttt{run\_demo.py} & 344 & CLI entrypoint. Orchestrates the entire pipeline: load data, build models, run experiments, generate figures. \\
\texttt{project\_config.py} & 207 & All hyperparameters and configuration in one place. Uses Python dataclasses. \\
\texttt{portfolio\_model.py} & 709 & Data loading (Excel and synthetic), sector aggregation, QUBO construction, Ising conversion. \\
\texttt{classical\_baselines.py} & 205 & Brute-force enumeration, greedy selection, continuous Markowitz, auto-penalty adjustment. \\
\texttt{qaoa\_simulator.py} & 224 & Pure-NumPy exact QAOA simulator for 8 qubits. No external quantum library needed. \\
\texttt{bloqade\_experiments.py} & 625 & Circuit construction, connectivity scheduling, noise modeling, full experiment loop. \\
\texttt{results\_metrics.py} & 179 & Per-bitstring evaluation (return, risk, VaR, CVaR) and distribution summaries. \\
\texttt{make\_figures.py} & 732 & All 10 presentation figures and CSV/JSON export utilities. \\
\midrule
\textbf{Total} & \textbf{3{,}225} & \\
\bottomrule
\end{longtable}
\end{center}

\subsection{Detailed File Explanations}

\subsubsection{\texttt{project\_config.py} --- Configuration}

This file defines four dataclass objects:

\begin{itemize}
    \item \texttt{ModelHyperparameters}: The 5 parameters that define a single QUBO model ($B$, $q$, $\lambda$, $\gamma$, $\eta$).
    \item \texttt{SweepConfig}: Ranges for hyperparameter sweeps when running the full grid.
    \item \texttt{ExecutionConfig}: Runtime settings (shots, optimizer restarts, which model versions and noise models to include).
    \item \texttt{RunConfig}: Top-level config combining all the above plus file paths.
\end{itemize}

It also defines \texttt{CANONICAL\_SECTORS} (the fixed 8-sector ordering that maps sectors to qubit indices), \texttt{SECTOR\_ALIASES} (fuzzy matching for workbook labels), and \texttt{COLUMN\_ALIASES} (fuzzy matching for column headers).

\subsubsection{\texttt{portfolio\_model.py} --- Data and QUBO}

\textbf{Data loading} (\texttt{load\_workbook\_dataset}):
\begin{itemize}
    \item Opens the Excel file with \texttt{openpyxl}.
    \item Auto-detects sheet names by keyword matching (``asset'', ``scenario'', ``cov'').
    \item Fuzzy-matches column headers to canonical names.
    \item Fuzzy-matches sector labels to the 8 canonical sectors.
    \item Falls back to synthetic data if the workbook is missing.
\end{itemize}

\textbf{Sector aggregation} (\texttt{aggregate\_to\_sectors}):
\begin{itemize}
    \item Builds the $8 \times 50$ aggregation matrix $A$.
    \item Computes $R_{\text{sec}} = R \cdot A^\top$.
    \item Derives $\bm{\mu}$, $\bm{\Sigma}$, capital, and liquidity vectors.
    \item Min--max normalizes capital and illiquidity.
\end{itemize}

\textbf{QUBO construction} (\texttt{build\_qubo\_model}):
\begin{itemize}
    \item Builds $Q$ using Equations~\ref{eq:qubo_diagonal} and~\ref{eq:qubo_offdiag}.
    \item Scales $Q$ by $\alpha = \max|Q_{ij}|$.
    \item Calls \texttt{qubo\_to\_ising} to derive $h$, $J$, and offset using Equations~\ref{eq:ising_h}--\ref{eq:ising_offset}.
\end{itemize}

\textbf{Ising conversion} (\texttt{qubo\_to\_ising}):
Implements the exact derivation from Section~5.3. Returns the field vector $\mathbf{h}$, coupling matrix $J$, and constant offset.

\subsubsection{\texttt{qaoa\_simulator.py} --- NumPy QAOA}

This is a self-contained, exact state-vector simulator:

\begin{itemize}
    \item \texttt{basis\_bit\_table}: Precomputes the $256 \times 8$ table mapping state indices to bit patterns.
    \item \texttt{precompute\_ising\_energies}: Evaluates $E(z)$ for all 256 basis states. This only needs to be done once per QUBO model.
    \item \texttt{simulate\_qaoa\_state}: The core simulation loop. Starts from $|+\rangle^{\otimes 8}$, alternates cost phases and mixer rotations.
    \item \texttt{apply\_rx\_layer}: Applies $R_X(2\beta)$ to each qubit using stride-based indexing on the state vector, avoiding the $256 \times 256$ matrix.
    \item \texttt{optimize\_qaoa\_angles}: Multi-start coordinate descent over the $2p$ parameters.
\end{itemize}

\subsubsection{\texttt{classical\_baselines.py} --- Ground Truth}

\begin{itemize}
    \item \texttt{exact\_bruteforce\_search}: Evaluates all 256 bitstrings.
    \item \texttt{greedy\_selection\_baseline}: Greedy sector-by-sector selection.
    \item \texttt{continuous\_markowitz\_baseline}: Projected gradient descent on continuous weights.
    \item \texttt{ensure\_feasible\_penalty}: Auto-doubles $\lambda$ until the QUBO global minimum is feasible.
\end{itemize}

\subsubsection{\texttt{bloqade\_experiments.py} --- Quantum Experiments}

This file manages the full experiment sweep:

\begin{itemize}
    \item For each (depth, seed, connectivity mode, noise model, noise strength, readout error) configuration:
    \begin{enumerate}
        \item Optimize QAOA angles (or reuse from same depth/seed).
        \item Build a Cirq circuit if Cirq is installed.
        \item Simulate the circuit (Cirq ideal or NumPy fallback).
        \item Apply noise (Bloqade native or approximate fallback).
        \item Sample bitstrings and evaluate portfolio metrics.
        \item Record summary statistics.
    \end{enumerate}
\end{itemize}

Key functions:
\begin{itemize}
    \item \texttt{interaction\_edges\_from\_ising}: Extracts weighted edges from $J$.
    \item \texttt{greedy\_edge\_coloring}: Partitions edges into parallel groups.
    \item \texttt{build\_cirq\_qaoa\_circuit}: Constructs the Cirq circuit with the CNOT/RZ/CNOT decomposition.
    \item \texttt{approximate\_noisy\_probabilities}: The mixing noise model.
    \item \texttt{simulate\_with\_noise\_if\_available}: Tries Bloqade noise first, falls back to approximate.
\end{itemize}

\subsubsection{\texttt{results\_metrics.py} --- Evaluation}

\begin{itemize}
    \item \texttt{evaluate\_selection}: Given a bitstring, computes all financial and QUBO metrics (return, variance, volatility, VaR, CVaR, capital, liquidity, QUBO objective).
    \item \texttt{summarize\_probability\_distribution}: Given a full QAOA probability vector, computes $P(\text{optimal})$, $P(\text{feasible})$, and expected objective.
    \item \texttt{sample\_counts\_from\_probabilities}: Simulates measurement shots by drawing from the probability distribution.
\end{itemize}

\subsubsection{\texttt{make\_figures.py} --- Visualization}

Generates all 10 presentation figures. Each function:
\begin{itemize}
    \item Takes processed data (records, aggregated dataset, or QUBO model).
    \item Creates a Matplotlib figure.
    \item Saves to PNG at 220 DPI.
\end{itemize}

Also provides \texttt{write\_csv} and \texttt{write\_json} helpers for result export, and \texttt{export\_qubo\_artifacts} for saving the $Q$, $J$ matrices and model metadata.

\subsubsection{\texttt{run\_demo.py} --- Pipeline Orchestrator}

The main function:
\begin{enumerate}
    \item Parses CLI arguments.
    \item Loads data (workbook or synthetic).
    \item Aggregates to sectors.
    \item For each (hyperparameter set, model version):
    \begin{enumerate}
        \item Builds the QUBO model.
        \item Validates the penalty (ensure feasible optimum).
        \item Exports QUBO artifacts.
        \item Runs classical baselines.
        \item Runs quantum experiments.
        \item Collects all results.
    \end{enumerate}
    \item Writes all CSVs and JSONs.
    \item Generates all 10 figures.
\end{enumerate}


% ══════════════════════════════════════════════════════════════════════════════
\section{Key Design Decisions and Why}
% ══════════════════════════════════════════════════════════════════════════════

\subsection{Why 8 Qubits?}

The Hartford workbook naturally groups 50 assets into 8 sectors. Rather than
artificially expanding or reducing the qubit count, we matched it to the
problem's natural structure. This means every qubit has a clear business
meaning: ``include this sector or not.''

\subsection{Why $B = 4$?}

Selecting half of 8 sectors gives $\binom{8}{4} = 70$ feasible portfolios---a
non-trivial search space. $B = 1$ or $B = 7$ would make the problem trivially
easy. $B = 4$ also mirrors common portfolio diversification practice.

\subsection{Why Equal Weights?}

Allowing continuous weights within selected sectors would turn the problem into
a mixed-integer program: discrete sector selection plus continuous weight
optimization. The continuous part would need to be handled separately from the
quantum optimization. Equal weighting keeps the problem purely binary, which
is exactly what QUBO/QAOA is designed for.

\subsection{Why Not Use SciPy for QAOA Optimization?}

We deliberately avoided SciPy to minimize dependencies and ensure the project
runs in any Python environment with just NumPy. The coordinate descent
optimizer is simple but sufficient for demonstrating the QAOA workflow. At
8 qubits, the improvement from a better optimizer would change $P(\text{optimal})$
by perhaps $0.1$--$0.5$ percentage points---not enough to change the conclusion
that QAOA is approximately uniform over feasible states.

\subsection{Why Four Model Variants?}

Running base, capital, liquidity, and both variants serves multiple purposes:
\begin{itemize}
    \item Shows that the QUBO formulation is flexible and can incorporate domain-specific constraints.
    \item Demonstrates to Hartford stakeholders how different regulatory priorities change the optimal portfolio.
    \item Validates that the quantum pipeline handles all four variants identically (same code path, same correctness).
\end{itemize}

\subsection{Why Include Brute-Force?}

At $n = 8$, brute-force is instant. Including it gives us an unambiguous ground
truth for every experiment. This is honest and makes all quantum claims
verifiable. It also lets us compute the exact energy landscape, which explains
\textit{why} QAOA performs the way it does.


% ══════════════════════════════════════════════════════════════════════════════
\section{What We Tried That Did Not Work}
% ══════════════════════════════════════════════════════════════════════════════

\subsection{H/CZ Decomposition for ZZ Gates}

Our initial Cirq circuit used the decomposition:
\[
e^{-i\theta Z_i Z_j} \stackrel{?}{=} H_j \cdot CZ(i,j) \cdot R_Z(2\theta)_j \cdot CZ(i,j) \cdot H_j
\]

This is \textbf{incorrect}. The $H$--$CZ$ sequence changes the basis to compute
$XX$ interactions, not $ZZ$. We discovered this when the Cirq circuit
probabilities did not match the NumPy simulator. The fix was switching to the
standard CNOT/RZ/CNOT decomposition (Equation~\ref{eq:zz_decomp}).

\textbf{Lesson:} Always validate quantum circuits against a known-correct
reference implementation. Our NumPy simulator served this role.

\subsection{Sharper QAOA Concentration}

We hoped that increasing QAOA depth from $p = 1$ to $p = 3$ would concentrate
probability mass on the optimum. It did not make a meaningful difference.

The reason is visible in the energy landscape: the gap between the best and
second-best feasible solutions is $\sim 0.000031$ in scaled units. For QAOA to
distinguish solutions this close in energy, it would need either much deeper
circuits (large $p$) or a fundamentally different approach like warm-starting.

\subsection{Pruned Connectivity}

We tried keeping only the top-12 strongest Ising couplings to reduce gate count.
While this reduced circuit depth, it also degraded solution quality by about
$50\%$ (from $\sim 0.009$ to $\sim 0.004$ optimal probability). Even weak
couplings carry information.

\textbf{Lesson:} Pruning is useful for depth reduction but must be balanced
against information loss. A smarter approach might prune edges below a threshold
based on the energy gap they contribute to.

\subsection{Running on Python 3.9}

The code uses \texttt{@dataclass(slots=True)}, which requires Python 3.10+.
We discovered this when trying to run on a machine with only Python 3.9. The
\texttt{slots=True} parameter was used for memory efficiency but is not
essential. It could be removed for broader compatibility.


% ══════════════════════════════════════════════════════════════════════════════
\section{Figures: What Each One Shows and Why}
% ══════════════════════════════════════════════════════════════════════════════

\subsection{Portfolio Tradeoff Frontier}
\texttt{portfolio\_tradeoff\_frontier.png}

Scatter plot of the four exact optimal portfolios in (volatility, expected return)
space. Point size encodes liquidity; color encodes capital charge.

\textbf{Purpose:} Shows judges immediately how the insurance constraints change
the optimal portfolio. Base model chases return (upper right); capital model
retreats to safety (lower left).

\subsection{Scenario Return Distributions}
\texttt{scenario\_return\_distributions.png}

Kernel density estimates of the 1{,}200-scenario return distributions for each
optimal portfolio. CVaR95 values are shown in the legend.

\textbf{Purpose:} Demonstrates that we actually use the scenario data, not just
summary statistics. Makes tail risk visible.

\subsection{Energy Landscape}
\texttt{energy\_landscape\_base.png}

All 256 bitstrings ranked by QUBO objective. Feasible states are highlighted.
The gap annotation shows why QAOA struggles.

\textbf{Purpose:} The single most important explanatory figure. Answers ``why
doesn't QAOA find the optimum with high probability?'' before the question is
asked.

\subsection{QAOA Parameter Landscape}
\texttt{qaoa\_parameter\_landscape\_base\_p1.png}

Heatmap of expected Ising energy as a function of $(\gamma, \beta)$ for $p = 1$.
The optimized point is marked.

\textbf{Purpose:} Shows that the QAOA landscape has structure (basins and ridges)
and that the optimizer finds a reasonable minimum.

\subsection{QAOA Convergence}
\texttt{qaoa\_convergence\_base.png}

Optimizer trace (expected energy vs.\ step) for $p = 1, 2, 3$.

\textbf{Purpose:} Shows the optimization process. The random restart phase
(high variance) followed by coordinate descent refinement (convergence) is
clearly visible.

\subsection{Classical vs Quantum Summary Table}
\texttt{classical\_vs\_quantum.png}

Comparison table showing exact bitstrings, greedy match, continuous gap,
quantum bitstrings, match status, and probabilities for all four models.

\textbf{Purpose:} The definitive comparison. All four quantum results match
exact, but the probabilities are $\sim 1\%$.

\subsection{Quantum Probability Summary}
\texttt{quantum\_probability\_summary.png}

Left panel: absolute $P(\text{optimal})$ and $P(\text{feasible})$.
Right panel: $P(\text{optimal} \mid \text{feasible})$ compared to uniform $1/70$.

\textbf{Purpose:} Makes the honest limitation explicit. The conditional
probability is indistinguishable from uniform.

\subsection{Atom Layout}
\texttt{atom\_layout.png}

2D MDS embedding of sectors based on correlation distance. Edge thickness
encodes correlation strength.

\textbf{Purpose:} Neutral-atom tie-in. Shows how sector correlations could
inform physical atom placement on hardware.

\subsection{Connectivity Tradeoff}
\texttt{connectivity\_tradeoff.png}

Scatter plot of average circuit depth vs.\ average $P(\text{optimal})$ for the
four connectivity modes.

\textbf{Purpose:} Shows that smart scheduling (hand/auto) achieves the same
quality at much lower depth. Pruning trades quality for depth.

\subsection{Success vs Noise}
\texttt{success\_vs\_noise.png}

$P(\text{optimal})$ as a function of noise strength for each noise model.
Noiseless baseline shown as dashed line.

\textbf{Purpose:} Quantifies noise sensitivity. Essential for the hardware
analysis portion of the presentation.


% ══════════════════════════════════════════════════════════════════════════════
\section{Glossary of Key Terms}
% ══════════════════════════════════════════════════════════════════════════════

\begin{description}[leftmargin=3cm, style=nextline]
    \item[QUBO] Quadratic Unconstrained Binary Optimization. An optimization
    problem over binary variables with a quadratic objective function.

    \item[Ising Model] A physical model of interacting spins. Equivalent to QUBO
    up to a linear change of variables.

    \item[QAOA] Quantum Approximate Optimization Algorithm. A variational quantum
    algorithm for combinatorial optimization.

    \item[Hamiltonian] In quantum mechanics, the operator whose eigenvalues are
    the system's energy levels. The cost Hamiltonian $H_C$ encodes our
    optimization objective.

    \item[Statevector] The full quantum state of the system. For $n$ qubits,
    a complex vector of dimension $2^n$.

    \item[Pauli-$Z$] The diagonal Pauli matrix $Z = \text{diag}(1, -1)$. Its
    eigenvalues $\pm 1$ correspond to the computational basis states $|0\rangle$
    and $|1\rangle$.

    \item[Pauli-$X$] The bit-flip matrix $X = \bigl(\begin{smallmatrix} 0 & 1 \\ 1 & 0 \end{smallmatrix}\bigr)$. The QAOA mixer Hamiltonian is a sum of $X$ operators.

    \item[$R_X(\theta)$] Rotation around the $X$-axis by angle $\theta$:
    $e^{-i\theta X/2}$.

    \item[$R_Z(\theta)$] Rotation around the $Z$-axis by angle $\theta$:
    $e^{-i\theta Z/2}$.

    \item[CNOT] Controlled-NOT gate. Flips the target qubit if the control qubit
    is $|1\rangle$.

    \item[Circuit Depth] The number of sequential time steps in a quantum circuit.
    Gates in the same time step are executed in parallel.

    \item[Chromatic Index] The minimum number of colors needed to edge-color a
    graph such that no two adjacent edges share a color.

    \item[VaR] Value at Risk. A quantile-based risk measure.

    \item[CVaR] Conditional Value at Risk (Expected Shortfall). The expected loss
    in the tail beyond VaR.

    \item[Markowitz] Harry Markowitz's mean--variance portfolio theory (1952).

    \item[Bloqade] QuEra's software framework for programming neutral-atom quantum
    computers.

    \item[Cirq] Google's Python framework for writing and simulating quantum
    circuits.
\end{description}


% ══════════════════════════════════════════════════════════════════════════════
\section{Anticipated Judge Questions}
% ══════════════════════════════════════════════════════════════════════════════

\begin{enumerate}
    \item \textbf{``Is there a quantum advantage here?''}

    No, and we are explicit about this. At 8 qubits, brute-force enumeration
    solves the problem exactly in microseconds. The value of this project is in
    the end-to-end pipeline, the hardware analysis, and the demonstration that
    the workflow scales to larger instances where classical methods fail.

    \item \textbf{``Why is QAOA basically uniform over feasible states?''}

    The energy gap between the best and second-best feasible solution is
    $\sim 3 \times 10^{-5}$ in scaled units. QAOA at shallow depth cannot
    resolve energy differences this small. The energy landscape figure
    (Section~10.3) shows this clearly.

    \item \textbf{``What would it take for quantum advantage?''}

    Scaling to 16--32 qubits (individual assets instead of sectors, or finer
    sector granularity) would make brute-force infeasible ($2^{16} = 65{,}536$
    to $2^{32} \approx 4 \times 10^9$ states). At that scale, QAOA or quantum
    annealing could become competitive, especially if the energy gaps are larger
    in the unsimplified problem.

    \item \textbf{``How realistic is your noise model?''}

    Our approximate noise model (uniform mixing) is a placeholder. It captures
    the qualitative behavior (more noise $\to$ worse performance) but is not
    calibrated to real hardware. When Bloqade is installed, the code upgrades
    to hardware-calibrated noise transforms automatically.

    \item \textbf{``Why not use a more sophisticated optimizer?''}

    Coordinate descent was chosen for simplicity and zero external dependencies.
    A better optimizer (COBYLA, Nelder--Mead) would find marginally better angles
    but would not change the conclusion: the problem's energy landscape is too
    flat for shallow QAOA to exploit.

    \item \textbf{``What does the neutral-atom layout actually tell you?''}

    The MDS layout places correlated sectors nearby. On a real neutral-atom
    device, atoms that need strong ZZ interactions should be physically close
    (interaction strength decays with distance). Our layout provides a starting
    point for mapping the problem onto hardware.

    \item \textbf{``Why did you choose equal weighting instead of optimizing weights?''}

    Equal weighting keeps the problem purely binary (QUBO-native). Continuous
    weight optimization would require a different quantum formulation, such as
    CVQE or a hybrid approach with classical inner-loop optimization.

    \item \textbf{``Can this handle more than 8 sectors?''}

    Yes. The QUBO construction, Ising conversion, and QAOA simulation are all
    parameterized by $n$. The only 8-specific piece is the brute-force baseline,
    which would need to be replaced by sampling-based evaluation at $n > 20$.
\end{enumerate}


% ══════════════════════════════════════════════════════════════════════════════
\section{Mathematical Reference Card}
% ══════════════════════════════════════════════════════════════════════════════

For quick reference during the presentation:

\begin{align}
\text{QUBO objective:} \quad & f(\mathbf{x}) = \mathbf{x}^\top Q\, \mathbf{x} \\[8pt]
\text{Full insurer objective:} \quad & f(\mathbf{x}) = \frac{q}{B^2}\, \mathbf{x}^\top \bm{\Sigma}\, \mathbf{x}
    - \frac{1}{B}\, \bm{\mu}^\top \mathbf{x}
    + \lambda (\mathbf{1}^\top \mathbf{x} - B)^2
    + \frac{\gamma}{B}\, \tilde{\mathbf{c}}^\top \mathbf{x}
    + \frac{\eta}{B}\, \tilde{\bm{\imath}}^\top \mathbf{x} \\[8pt]
\text{QUBO diagonal:} \quad & Q_{ii} = \frac{q}{B^2}\Sigma_{ii} - \frac{\mu_i}{B} + \lambda(1-2B) + \frac{\gamma}{B}\tilde{c}_i + \frac{\eta}{B}\tilde{\imath}_i \\[8pt]
\text{QUBO off-diag:} \quad & Q_{ij} = \frac{2q}{B^2}\Sigma_{ij} + 2\lambda \quad (i < j) \\[8pt]
\text{Bit to spin:} \quad & x_i = \frac{1 - s_i}{2} \\[8pt]
\text{Ising Hamiltonian:} \quad & H_C = \sum_i h_i Z_i + \sum_{i<j} J_{ij} Z_i Z_j \\[8pt]
\text{QAOA state:} \quad & |\bm{\gamma},\bm{\beta}\rangle = \prod_{k=1}^p e^{-i\beta_k H_M} e^{-i\gamma_k H_C} |+\rangle^{\otimes n} \\[8pt]
\text{ZZ gate:} \quad & e^{-i\theta Z_i Z_j} = \text{CNOT}(i,j) \cdot R_Z(2\theta)_j \cdot \text{CNOT}(i,j) \\[8pt]
\text{Feasible count:} \quad & \binom{8}{4} = 70 \\[8pt]
\text{VaR}_\alpha: \quad & \inf\{l : P(\text{Loss} > l) \leq 1 - \alpha\} \\[8pt]
\text{CVaR}_\alpha: \quad & \mathbb{E}[\text{Loss} \mid \text{Loss} \geq \text{VaR}_\alpha]
\end{align}


\end{document}
